\documentclass[11pt]{article}
\usepackage{fullpage, fancyhdr, ifthen, amssymb, amsfonts, amsmath, booktabs}
\usepackage{tabularx}
\usepackage{multicol}
\usepackage[margin=0.25in]{geometry}
\usepackage{upquote}
\usepackage{forest}
\usepackage{ulem}
\usepackage{colortbl}
\usepackage{enumitem}
\pagestyle{empty}

\begin{document}

\begin{table}[h]
\centering
\arrayrulecolor{black!10} % Very light table outline
\begin{tabular}{|p{\textwidth}|}
\hline
\begin{minipage}[t]{\textwidth}
\vspace{0pt}  % This is crucial for top alignment
\noindent \textbf{Instructions:} Please answer the following questions, making sure to write legibly on the following pages. The code below is a subsection of a Python file that provides context for the questions below.
\begin{itemize}
\item Since we are writing Python code, please be careful about spacing/indentation. Leave enough room that it is obvious when you want things indented. 
\item Make sure to write legibly.
\item All questions are worth equal points unless otherwise stated.
\item The quiz is cumulative and may include material from previous weeks.
\end{itemize}
\end{minipage} 
\\
\hline
\end{tabular}
\end{table}

{\color{lightgray}\hrule}
\begin{verbatim}
from flask import Flask, request, jsonify
from logger_utils import logger  # Assume this logger is already imported

app = Flask(__name__)

def calculate_discount(price, discount_percent):
    # Calculate the final price after applying a discount.    

    return price * (1 - discount_percent / 100)

def apply_multiple_discounts(price, *discounts):
    # Apply multiple discounts sequentially to a price.

    result = price
    for discount in discounts:
        result = calculate_discount(result, discount)
    return result

@app.route('/api/products/<product_id>/price', methods=['GET'])
def get_product_price(product_id):
    discount = request.args.get('discount')
    base_price = get_base_price(product_id)  # Assume this function exists
    final_price = calculate_discount(base_price, discount)
    return jsonify({
        'product_id': product_id,
        'base_price': base_price,
        'discount_percent': discount,
        'final_price': round(final_price, 2)
    }), 200

\end{verbatim}
{\color{lightgray}\hrule}

\clearpage

\begin{tabularx}{\textwidth}{l|X}
\textbf{Quiz 7A} &   \textbf{Name: } \\
\hline
\end{tabularx}


\begin{enumerate}

\item (3 Points) What does CRUD stand for in the context of API development? Please just write the words, no explanation needed.
\vspace{1.5cm}


\item Write a single test function called \texttt{test\_calculate\_discount} that tests the \texttt{calculate\_discount} function shown on the first page. Assume a logger named \texttt{logger} is already imported (as shown in the code context). It should take zero arguments and have no decorators. Your test should:
\begin{itemize}
    \item Test that a price of 100 with a 10\% discount returns 90.0
    \item Test that a price of 50 with a 20\% discount returns 40.0
    \item Log each successful test result at the INFO level using the format: \\ \texttt{"Test passed: \{price\}, \{discount\}, \{result\}"}
\end{itemize}

\vspace{5.5cm}
	


\item Write a single integration test function called \texttt{test\_discount\_edge\_cases} that tests \texttt{apply\_multiple\_discounts} (and so explicitly tests how  the functions work together.). Your test should:
\begin{itemize}
    \item Test that applying zero discounts (empty list) returns the original price
    \item Test that applying a 0\% discount produces the same result as applying no discount
    \item Test that applying a 100\% discount results in a final price of 0.0
    \item Log each successful test result using the same format as in Question 2.
\end{itemize}

\vspace{6.5cm}

\item (2 Points) In class we discussed autodocs systems like MkDocs. Please explain why the \texttt{site/} directory should {\it not} be included in a git repository. Write a single sentence answer.


\end{enumerate}

\end{document}

